\documentclass[12pt,letterpaper]{article}

% --- Codificación, Fuentes e Idioma (Motor PDFLaTeX) ---
\usepackage[utf8]{inputenc}
\usepackage[T1]{fontenc}
\usepackage[spanish,mexico]{babel}
\usepackage{newtxtext,newtxmath} % Times New Roman de alta calidad para APA 7 (Texto y Matematicas)

% --- Geometría y Márgenes (APA 7: Exactamente 1 pulgada / 2.54 cm) ---
\usepackage[margin=1in]{geometry}

% --- Espaciado y Párrafos ---
\usepackage{setspace}
\onehalfspacing % Interlineado 1.5 común en reportes universitarios de ingeniería
\usepackage{parskip} % Espacio entre párrafos en lugar de sangría forzada excesiva

% --- Matemáticas y Símbolos Avanzados ---
\usepackage{amsmath,amssymb,amsfonts,amsthm}
\usepackage{icomma} % Corrige el comportamiento de la coma decimal en español

% --- Gráficos, Tablas y Entornos Flotantes ---
\usepackage{graphicx}
\usepackage{float}
\usepackage{booktabs} % Tablas profesionales (reemplaza líneas verticales feas por horizontales limpias)
\usepackage{caption}
\usepackage{subcaption}

% --- Enlaces y Referencias cruzadas inteligentes ---
\usepackage{hyperref}
\hypersetup{
    colorlinks=true,
    linkcolor=black,
    citecolor=black,
    urlcolor=blue,
    pdftitle={Documento Tecnico},
    pdfauthor={Daniel}
}

% --- Bibliografía APA 7 con BibLaTeX ---
\usepackage[backend=biber,style=apa,sortcites=true,sorting=nyt,natbib=true]{biblatex}
\addbibresource{biblio/referencias.bib}

% --- Inserción de Código Limpio y Estilizado (Ingeniería en Sistemas) ---
\usepackage{listings}
\usepackage{xcolor}

\definecolor{bgcode}{rgb}{0.96,0.96,0.96}
\definecolor{keywordcolor}{rgb}{0.13,0.31,0.58}
\definecolor{stringcolor}{rgb}{0.06,0.55,0.06}
\definecolor{commentcolor}{rgb}{0.38,0.38,0.38}

\lstdefinestyle{sistemasstyle}{
    backgroundcolor=\color{bgcode},
    basicstyle=\ttfamily\small,
    breakatwhitespace=false,
    breaklines=true,
    captionpos=b,
    commentstyle=\color{commentcolor}\itshape,
    keywordstyle=\color{keywordcolor}\bfseries,
    stringstyle=\color{stringcolor},
    numberstyle=\tiny\color{gray},
    numbers=left,
    numbersep=8pt,
    showspaces=false,
    showstringspaces=false,
    showtabs=false,
    tabsize=4,
    frame=single,
    rulecolor=\color{lightgray}
}
\lstset{style=sistemasstyle}

% --- Encabezados y Pies de Página Profesionales ---
\usepackage{fancyhdr}
\pagestyle{fancy}
\fancyhf{}
\fancyhead[L]{\small Ingeniería en Sistemas Computacionales}
\fancyhead[R]{\small\thepage}
\renewcommand{\headrulewidth}{0.4pt}

% =================================================================
% INICIO DEL DOCUMENTO
% =================================================================
\begin{document}

% --- Portada Académica Formal ---
\begin{center}
    {\Large \textbf{INSTITUTO TECNOLÓGICO / UNIVERSIDAD}} \\[0.3cm]
    {\large Facultad de Ingeniería en Sistemas Computacionales} \\[0.2cm]
    {\normalsize 5to Semestre} \\[1.5cm]
    
    {\LARGE \textbf{Título Principal del Proyecto o Práctica}} \\[0.5cm]
    {\large Subtítulo opcional o descripción del módulo} \\[2cm]
    
    \begin{tabular}{rl}
        \textbf{Alumno:} & Daniel \\
        \textbf{Materia:} & Nombre de la Asignatura \\
        \textbf{Profesor:} & Nombre del Docente \\
        \textbf{Fecha:} & \today \\
    \end{tabular}
\end{center}

\thispagestyle{empty}
\newpage

% --- Sección 1: Introducción ---
\section{Introducción}
\label{sec:introduccion}

La redacción de este documento sigue los lineamientos del formato APA 7 adaptados para reportes de ingeniería. Los párrafos se estructuran de manera lógica y fluida. Es posible realizar referencias bibliográficas automáticas utilizando BibLaTeX, por ejemplo, al estudiar la ingeniería de software moderna \cite{sommerville2019}.

% --- Sección 2: Marco Teórico ---
\section{Marco Teórico}
\label{sec:marco_teorico}

\subsection{Fundamentos de la Arquitectura}
Desarrollo conceptual de los temas correspondientes a la práctica o investigación del semestre.

% --- Sección 3: Desarrollo e Implementación ---
\section{Desarrollo e Implementación}
\label{sec:desarrollo}

En esta sección se detallan los artefactos técnicos, diagramas y fragmentos de código desarrollados.

\subsection{Implementación de Código}
El Listado \ref{code:ejemplo} muestra un bloque de código estructurado en Python integrado directamente desde la carpeta de fuentes.

\begin{lstlisting}[language=Python, caption={Script de automatización base.}, label={code:ejemplo}]
def inicializar_sistema(puerto: int) -> bool:
    print(f"Iniciando servicio en el puerto {puerto}...")
    return True

if __name__ == "__main__":
    inicializar_sistema(8080)
\end{lstlisting}

\subsection{Diagramas del Sistema}
Para integrar esquemas, modelos Entidad-Relación o arquitecturas de red, utiliza el entorno gráfico estándar (Figura \ref{fig:arquitectura}).

\begin{figure}[H]
    \centering
    % \includegraphics[width=0.75\textwidth]{src/arquitectura.png}
    \caption{Arquitectura general del despliegue.}
    \label{fig:arquitectura}
\end{figure}

% --- Sección 4: Resultados y Discusión ---
\section{Resultados y Discusión}
\label{sec:resultados}
Análisis de rendimiento, pruebas unitarias o métricas obtenidas durante el desarrollo del proyecto.

% --- Sección 5: Conclusiones ---
\section{Conclusiones}
\label{sec:conclusiones}
Reflexiones finales sobre el cumplimiento de los objetivos planteados en la práctica de ingeniería.

% --- Referencias APA 7 ---
\newpage
\printbibliography[title={Referencias}]

\end{document}
